%% annex/p07/proof-complexity.tex --- proof-complexity sketches for Part VII

\section{The pigeonhole bottleneck argument}

\begin{annexitem}{annex:p7.complexity.1}{thm:p07-haken-resolution-lower-bound}{Haken's resolution lower bound}
\sketchnote{Haken's full bottleneck lemma with its numerical constants; the
count of critical partial injections falsifying each possible bottleneck
clause; and the bookkeeping that assigns every critical injection to a first
bottleneck node in a shared derivation DAG.  These are the combinatorial core
of the exponential estimate, not consequences of resolution soundness alone.}
For the exact formula \(\mathsf{PHP}^{n}_{n+1}\), call an assignment
\emph{critical} when it injects exactly \(n\) of the \(n+1\) pigeons into all
\(n\) holes and sets every variable belonging to the remaining pigeon to
false.  There are \((n+1)n!\) such assignments.  Each satisfies every
collision clause and falsifies exactly the totality clause belonging to its
unassigned pigeon.  The qualifier is necessary to distinguish this counted
family from arbitrary truth assignments.  It names a metalinguistic class of
assignments, not a new proof object; its unique-falsified-input property is the
nonredundant feature used by the routing and counting steps below.

Orient a resolution refutation from the empty clause toward its input clauses.
Whenever a critical assignment falsifies a resolvent, its value on the pivot
selects a parent clause that it also falsifies.  Following these choices sends
every critical assignment from the empty clause to the unique input totality
clause that it falsifies.  DAG sharing does not change this routing; it is the
reason assignments, rather than proof paths, are counted.

The load-bearing combinatorial input is Haken's bottleneck lemma in the
formulation surveyed by Beame and Pitassi.  Its clause-complexity predicate is
the source's predicate on clauses and compatible partial injections; it is
strictly stronger than merely counting how many pigeon indices a clause
mentions.  The lemma supplies absolute constants \(\alpha,\beta>0\) and proves
both facts needed here: every routed critical assignment has a first
bottleneck occurrence, and any fixed bottleneck occurrence receives at most a
\(2^{-\beta n}\) fraction of all critical assignments.  The proof of those two
facts is exactly the omitted bottleneck/counting argument named in the opening
sketch note; no independent-literal heuristic is substituted for it.

Charge each critical assignment to that first bottleneck occurrence.  Since
all \((n+1)n!\) assignments must be charged and one occurrence receives at
most a \(2^{-\beta n}\) fraction of them, the refutation contains at least
\(2^{\beta n}\) bottleneck clause occurrences.  Taking \(c=2^{\beta}>1\)
yields the theorem's bound.  The cited sources prove the imported lemma for
general DAG-like resolution and for precisely the nonfunctional CNF fixed in
Definition~\ref{def:p07-pigeonhole-cnf}
\cite{Haken1985,BeamePitassi1998}.
\end{annexitem}
